% Generated by scripts/make_paper_tables.py from the frozen results.
% Do not edit by hand; the next build overwrites it.
\begin{table}[tbp]
\centering
\caption{The funding spread that would reconcile model and disclosed value, under two assumptions about the call}
\label{tab:spreads}
\small
\begin{tabular}{l r r r r r}
\toprule
Note & Disclosed & Model & Spread (bp) & Model, no call & Spread (bp) \\
\midrule
N1 & 938.90 & 910.98 & -123 & n/a & -1 \\
N2 & 943.10 & 894.38 & -204 & n/a & -69 \\
N3 & 983.40 & 962.56 & -179 & n/a & +76 \\
N4 & 999.00 & 979.84 & no solution & n/a & +99 \\
N5 & 972.10 & 960.69 & -278 & n/a & n/a \\
N6 & 967.30 & 954.01 & -156 & n/a & +121 \\
\bottomrule
\end{tabular}
\notetext{\textit{Notes:} Each spread is the constant addition to the Treasury discount curve that sets the model value equal to the disclosed value, solved on a common set of paths. A funding spread should be positive and wider than the issuer's own senior bonds: the A-rated US corporate option-adjusted spread was 69\,bp on the pricing date. Under optimal exercise every solved spread is negative, from -278\,bp to -123\,bp, which would require the issuer to fund below Treasuries. \textit{No solution} means the disclosed value lies outside the bracket searched; \S\ref{sec:results} reports where that note's crossing lies. The auto-callable note has no no-call column: its redemption is contractual, so removing it describes a different security rather than a different modelling assumption.}
\end{table}
